\gddchapter{Testing sequence}
The prototype can be played with both voice and keyboard and the navigation modes can be swapped with a button press. 


\section{Novel version B - Voice navigation}

Since the feedback from the last 3 tests the testing is now done in reverse order - the participant is playing the voice version first and then switching to the keyboard baseline version.


\subsection{Observation log}
\begin{table}[H]
\centering
\begin{tabular}{|p{5cm}|p{10cm}|}
\hline
\textbf{Field} & \textbf{Response} \\ \hline
\textbf{Physical comfort} How is the participant feeling & Apart from the text being very hard to read the participant looks comfortable.\\[1.5cm] \hline
\textbf{Natural language observation} (is the player using long sentences or simple commands? What type of language do they use? Are they speaking first or 3rd person?) & The participant is using very complex commands using full blown natural language. \\[1.5cm] \hline
\textbf{Linguistic guessing} (is the particpant trying to guess the possible commands by speaking to the game?) & The player is trying to guess what's possible by speaking to the game. He also expressed the want to ask game for answers. \\[1.5cm] \hline
\textbf{Processing time user reaction} (Does the wait time break the suspension of disbelief) & It doesn't seem to be affecting this parcicipant. \\[1.5cm] \hline
\textbf{Voice processing pipeline edge cases} (System edge cases eg. two vs to)& There were no huge errors detected in his experience. \\[1.5cm] \hline
\end{tabular}
\end{table}


\section{Baseline A - Keyboard navigation}
After playing the audio version of the game, the player is now experiencing the baseline version with keyboard usage.
This allows the player to anchor in the typical experience and serves as a comparative point in the A/B testing.

\subsection{Observation log}
\begin{table}[h]
\centering
\begin{tabular}{|p{5cm}|p{10cm}|}
\hline
\textbf{Field} & \textbf{Response} \\ \hline
\textbf{Physical comfort} How is the participant feeling & The participant looks just as comfortable as in audio version.\\[1.5cm] \hline
\textbf{Immersion markers} (flow \& frustration) & The participant is struggling a little bit with immersion due to the fact that he's in need of new glasses and he's struggling with reading the descriptions.\\[1.5cm] \hline
\end{tabular}
\end{table}


\subsection{Sticky moments}
The moments that stood out during testing marked with timestamps.

\begin{table}[h]
\centering
\begin{tabular}{|p{3cm}|p{10cm}|}
\hline
\textbf{Timestamp} & \textbf{Log} \\ \hline
06:47 & Start of the playtest \\ \hline 
07:20 & Text too small \\ \hline 
08:55 & Can I ask questions? \\ \hline 
18:20 & Text too small again \\ \hline 
\end{tabular}
\end{table}

\section{Alignment check}
How much was the researcher participating in the gameplay experience? 

The researcher did his best not to participate actively during the testing sequence. 



